\section{The wire}
\label{app:wire}

A request to \op{POST /v1/decisions}, in OpenRouter's Decisions schema:

\begin{lstlisting}[style=wire]
{
  "model": "kai-1",
  "state": "Hi, we were billed twice for March. Please refund the duplicate today or we will cancel our plan.",
  "questions": {
    "department": {
      "type": "choice",
      "instructions": "Which department should handle this?",
      "criteria": {
        "billing":   "invoices, payments, refunds",
        "technical": "bugs, outages, system errors",
        "other":     "everything else"
      }
    },
    "urgency": {
      "type": "score",
      "instructions": "How urgent is this?",
      "criteria": ["not urgent", "soon", "blocking"]
    },
    "churn_risk": {
      "type": "noul",
      "instructions": "Does the user threaten to cancel or leave?"
    }
  }
}
\end{lstlisting}

The response carries one typed answer per question. The probabilities below
are those the English checkpoint gave in the recorded run; \op{confidence} is
Equation~\eqref{eq:kappa}, \op{answer\_confidence} is $\pmax$, and the
remaining Hanzo fields are omitted. The \op{id} and \op{usage} values are
illustrative.

\begin{lstlisting}[style=wire]
{
  "id": "dec_...",
  "model": "kai-1",
  "provider": "hanzo",
  "answers": {
    "department": {
      "type": "choice", "choice": "billing",
      "probabilities": {"billing": 0.9865, "technical": 0.008, "other": 0.0055},
      "confidence": 0.9798, "answer_confidence": 0.9865
    },
    "urgency": {
      "type": "score", "score": 1.7722,
      "legend": {"0": "not urgent", "1": "soon", "2": "blocking"},
      "probabilities": {"0": 0.0389, "1": 0.15, "2": 0.8111},
      "confidence": 0.7167, "answer_confidence": 0.8111
    },
    "churn_risk": {
      "type": "noul", "noul": 0.879,
      "confidence": 0.758, "answer_confidence": 0.879
    }
  },
  "usage": {"input_tokens": 164, "output_tokens": 0, "cost": 0}
}
\end{lstlisting}

Changing \op{"model"} to \op{"typesafe/jev-1.13"} sends the same request to
Jev through OpenRouter; nothing else in the request changes.

\section{A Decision Program}
\label{app:program}

\op{agent.command-risk@1} as defined for the Rust runtime \proposed:

\begin{lstlisting}[style=wire]
{
  "id": "agent.command-risk@1",
  "description": "Before an agent runs a shell command: allow it, ask the user, or deny it.",
  "model": "kai-1-agent",
  "mode": "advisory",
  "calibration": "cal_6dbb072f8109a222",
  "questions": {
    "verdict": {
      "type": "choice",
      "instructions": "Should the agent run this shell command?",
      "criteria": {
        "allow": "read-only, or confined to the working repository",
        "ask":   "writes outside the repository, installs software or touches the network",
        "deny":  "reads or sends secrets, keys, ~/.ssh or .env files, or destroys data"
      }
    },
    "secret_exposure":   {"type": "noul", "instructions": "Could this command read, print or send credentials, keys or tokens?"},
    "production_impact": {"type": "noul", "instructions": "Could this command change or break a production system or shared data?"},
    "network":           {"type": "noul", "instructions": "Does this command send or fetch data over the network?"}
  },
  "thresholds": {"verdict": 0.5, "secret_exposure": 0.3, "production_impact": 0.4, "network": 0.5},
  "verdict": {
    "from": "verdict",
    "escalate": {"secret_exposure": "deny", "production_impact": "ask", "network": "ask"}
  }
}
\end{lstlisting}

The verdict question abstains to \ask{} below $\kappa = 0.5$; each escalation
raises the verdict when its \op{noul} reaches its threshold; the result is
joined with the policy's verdict when the program runs enforced.
